\documentclass[11pt,a4paper]{article}

\usepackage[utf8]{inputenc}
\usepackage[T1]{fontenc}
\usepackage[spanish,provide=*]{babel}
\usepackage{csquotes}
\usepackage{charter}
\usepackage[scaled=0.92]{helvet}
\usepackage[a4paper,top=2.5cm,bottom=2.5cm,left=2.2cm,right=2.2cm,headheight=14pt]{geometry}
\usepackage{booktabs}
\usepackage{longtable}
\usepackage{array}
\usepackage{xcolor}
\usepackage{colortbl}
\usepackage{titlesec}
\usepackage{fancyhdr}
\usepackage{enumitem}
\usepackage{tcolorbox}
\usepackage{listings}
\usepackage{amssymb}
\usepackage{amsmath}
\usepackage{microtype}
\usepackage[hidelinks]{hyperref}

\definecolor{uteqBlue}{HTML}{0E3F7E}
\definecolor{uteqMid}{HTML}{1E5AA8}
\definecolor{uteqTeal}{HTML}{2E86A1}
\definecolor{uteqGreen}{HTML}{4FAE60}
\definecolor{uteqAmber}{HTML}{F2A413}
\definecolor{uteqSoft}{HTML}{EAF1F9}
\definecolor{uteqRed}{HTML}{B02418}
\definecolor{uteqGray}{HTML}{3A3A3A}
\definecolor{grisCodigo}{HTML}{F2F4F6}
\definecolor{filaReuso}{HTML}{EDF7EE}
\definecolor{filaCambia}{HTML}{FDF4E3}
\definecolor{filaNuevo}{HTML}{FBEDEC}

\hypersetup{colorlinks=true,linkcolor=uteqBlue,urlcolor=uteqMid,citecolor=uteqBlue}

\titleformat{\section}{\normalfont\Large\bfseries\color{uteqBlue}}{\thesection}{0.7em}{}
\titleformat{\subsection}{\normalfont\large\bfseries\color{uteqMid}}{\thesubsection}{0.6em}{}
\titleformat{\subsubsection}{\normalfont\normalsize\bfseries\color{uteqTeal}}{\thesubsubsection}{0.5em}{}
\titlespacing*{\section}{0pt}{16pt}{7pt}
\titlespacing*{\subsection}{0pt}{12pt}{5pt}

\pagestyle{fancy}
\fancyhf{}
\renewcommand{\headrulewidth}{0.4pt}
\fancyhead[L]{\footnotesize\color{uteqBlue}PFC Entrega 4 $\cdot$ Guía de reutilización $\cdot$ Equipo AGLS}
\fancyhead[R]{\footnotesize\color{uteqBlue}Aplicaciones Distribuidas (ISR-701) $\cdot$ \thepage}
\fancyfoot[C]{\footnotesize\color{uteqGray}Prof. Ing. Gleiston C. Guerrero-Ulloa, Mgs. $\cdot$ UTEQ}

\setlength{\parskip}{5pt}
\setlength{\parindent}{0pt}
\renewcommand{\arraystretch}{1.15}

\lstset{
	basicstyle=\ttfamily\scriptsize,
	backgroundcolor=\color{grisCodigo},
	frame=single,
	rulecolor=\color{uteqBlue},
	breaklines=true,
	columns=fullflexible,
	xleftmargin=5pt,
	xrightmargin=5pt,
	framexleftmargin=5pt
}

\newtcolorbox{uteqbox}[1]{colback=uteqSoft,colframe=uteqBlue,boxrule=0.8pt,arc=2pt,
	left=7pt,right=7pt,top=6pt,bottom=6pt,title={\bfseries #1},
	fonttitle=\normalsize\color{white},colbacktitle=uteqBlue}
\newtcolorbox{pisoBox}[1]{colback=white,colframe=uteqRed,boxrule=0.9pt,arc=2pt,
	left=7pt,right=7pt,top=6pt,bottom=6pt,title={\bfseries #1},
	fonttitle=\normalsize\color{white},colbacktitle=uteqRed}
\newtcolorbox{avisoBox}[1]{colback=white,colframe=uteqAmber,boxrule=0.9pt,arc=2pt,
	left=7pt,right=7pt,top=6pt,bottom=6pt,title={\bfseries #1},
	fonttitle=\normalsize\color{uteqGray},colbacktitle=uteqAmber}
\newtcolorbox{defBox}[1]{colback=white,colframe=uteqGreen,boxrule=0.9pt,arc=2pt,
	left=7pt,right=7pt,top=6pt,bottom=6pt,title={\bfseries #1},
	fonttitle=\normalsize\color{white},colbacktitle=uteqGreen}

\newcommand{\reuso}{\cellcolor{filaReuso}\textbf{REUTILIZA}}
\newcommand{\cambia}{\cellcolor{filaCambia}\textbf{MODIFICA}}
\newcommand{\nuevo}{\cellcolor{filaNuevo}\textbf{NUEVO}}

\begin{document}

\begin{titlepage}
	\centering
	\vspace*{0.8cm}
	{\small\color{uteqGray} UNIVERSIDAD TÉCNICA ESTATAL DE QUEVEDO\\ Facultad de Ciencias de la Computación\\ Carrera de Ingeniería de Software\par}
	\vspace{14pt}
	{\color{uteqBlue}\rule{\textwidth}{2pt}}\\[7pt]
	\colorbox{uteqBlue}{\parbox{0.96\textwidth}{\centering\vspace{5pt}{\large\bfseries\color{white} GUÍA DE REUTILIZACIÓN Y ENTREGA --- PFC ENTREGA 4}\vspace{5pt}}}\\[7pt]
	{\color{uteqBlue}\rule{\textwidth}{2pt}}\\[16pt]
	{\LARGE\bfseries\color{uteqBlue} Equipo AGLS --- TiendaTech\par}
	\vspace{5pt}
	{\large\color{uteqGray} Qué ya tienen, cómo se entrega ahora, qué cambia y qué es nuevo\par}
	\vspace{20pt}
	\begin{minipage}{0.93\textwidth}
		\small
		\begin{tabular}{@{}p{4.6cm}p{9.8cm}@{}}
			\toprule
			\textbf{Documento rector} & \emph{Guía Integral de la Entrega 4 del Proyecto Fin de Curso}, v1.0. Esta guía \textbf{no la sustituye}: se acopla a ella \\
			\textbf{Asignatura} & Aplicaciones Distribuidas (ISR-701), séptimo semestre \\
			\textbf{Período} & Regular 2026--2027 \\
			\textbf{Equipo} & AGLS (3--4 integrantes) \\
			\textbf{Semana de entrega} & \textbf{18} \\
			\textbf{Cierre} & \textbf{Viernes 4 de septiembre de 2026, 23:55} (último \emph{commit} admitido) \\
			\textbf{Fecha de esta guía} & Semana 17 --- quedan \textbf{nueve días naturales} hasta el cierre \\
			\textbf{Rama de trabajo} & \texttt{feature/entrega-4}, fusionada a \texttt{main} por \emph{pull request} \\
			\textbf{Calificación} & Según la rúbrica de nueve dimensiones de la guía rectora \\
			\bottomrule
		\end{tabular}
	\end{minipage}
	\vfill
	{\small\color{uteqGray} Quevedo, Ecuador\par}
\end{titlepage}

\tableofcontents
\newpage

\section{Cómo leer esta guía}

La Entrega 4 no empieza de cero. La guía rectora lo dice con claridad: la E4 hereda el análisis del dominio y el prototipo de comunicación de la E1, la arquitectura de microservicios REST con Docker Compose de la E2, y la capa de datos distribuida en CockroachDB con el \emph{pipeline} PySpark de la E3. \textbf{Ninguna de esas piezas se descarta.} La E4 las envuelve, las refactoriza y las expone a los usuarios finales por dos canales nuevos.

Esta guía traduce ese principio a la situación concreta del equipo AGLS. Cada bloque de trabajo aparece clasificado con una de tres etiquetas:

\begin{center}
	\small
	\begin{tabular}{@{}p{2.6cm}p{12.6cm}@{}}
		\toprule
		\reuso & El artefacto ya existe y se conserva tal como está. Solo hay que verificar que sigue funcionando y volverlo a presentar en el paquete final \\
		\addlinespace
		\cambia & El artefacto ya existe pero debe transformarse. Se indica exactamente qué se toca y qué se conserva \\
		\addlinespace
		\nuevo & No existe. Hay que construirlo desde cero en esta entrega \\
		\bottomrule
	\end{tabular}
\end{center}

\begin{uteqbox}{La regla que ordena todo el trabajo}
	Antes de escribir una línea de código nueva, el equipo debe hacer el inventario de la sección 2. La mayor parte de la Entrega 4 no es construcción: es \textbf{recuperar, verificar, refactorizar y volver a presentar}. Lo verdaderamente nuevo son las dos aplicaciones cliente y el instrumental de medición.
\end{uteqbox}

\begin{pisoBox}{Dónde está el equipo en el calendario}
	Esta guía se entrega en la \textbf{semana 17}. El cierre es el \textbf{viernes 4 de septiembre de 2026 a las 23:55}, hora del último \emph{commit} admitido: lo que se envíe después no se evalúa. Quedan \textbf{nueve días naturales}, dos de ellos fin de semana.\\[5pt]
	Por tanto, cuando esta guía menciona el trabajo de semanas anteriores no está proponiendo un plan: está describiendo lo que el equipo \textbf{debía tener ya}. Si algo de eso falta, no se planifica, se triaja: vaya directamente a la sección 10.
\end{pisoBox}

\section{Inventario: lo que AGLS ya tiene}

Esta tabla es el punto de partida y es una \textbf{auditoría, no una planificación}: todo lo que aparece aquí debía estar terminado en las entregas de las semanas 4, 7 y 11. El equipo debe recorrerla hoy mismo, marcar el estado real de cada artefacto y anotar dónde está en el repositorio. Un artefacto que no se encuentre hoy no es trabajo pendiente: es un hueco que hay que tapar en nueve días o asumir como pérdida consciente.

\begin{center}
	\footnotesize
	\begin{longtable}{@{}p{3.7cm}p{1.5cm}p{5.2cm}p{4.4cm}@{}}
		\toprule
		\textbf{Artefacto existente} & \textbf{Origen} & \textbf{Cómo se entrega ahora en la E4} & \textbf{Estado} \\
		\midrule
		\endhead
		Problema del dominio cuantificado y modelo de fallos & E1 & Se integra en la Introducción del manuscrito, con las cifras verificadas de nuevo & \cambia \\
		\addlinespace
		Diagramas C4 niveles 1 y 2 & E1 & Van a la sección de Arquitectura. El nivel 3 se rehace porque cambia (ver más abajo) & \reuso \\
		\addlinespace
		ADR-001 a ADR-004 & E1--E3 & Se conservan íntegros en \texttt{docs/adr/}. Se les añade el impacto que la E4 tuvo sobre ellos & \cambia \\
		\addlinespace
		Servidor y clientes de sockets TCP, servicio gRPC, relojes de Lamport & E1 \newline PE-U1 & Se conservan como el canal de reserva atómica de stock entre Órdenes e Inventario y como el orden total de reservas concurrentes & \reuso \\
		\addlinespace
		Microservicios REST con base de datos propia & E2 & Se refactorizan en cuatro capas (Módulo A). \textbf{La API externa no cambia} & \cambia \\
		\addlinespace
		API Gateway con enrutamiento, autenticación y límite de tasa & E2 & Se conserva. Ahora sirve a dos clientes en lugar de a ninguno & \reuso \\
		\addlinespace
		Especificación OpenAPI 3.0 & E2 & Se conserva y se verifica que sigue describiendo la API real & \cambia \\
		\addlinespace
		Autenticación con JWT y \texttt{/auth/login} & E2 & Es el mismo \emph{endpoint} que consumirán la web y la móvil & \reuso \\
		\addlinespace
		\texttt{docker-compose.yml} y \texttt{Dockerfile} multi-etapa & E2 & Se amplía con web, colector de trazas y servicios nuevos & \cambia \\
		\addlinespace
		Clúster CockroachDB de tres nodos y fragmentación del catálogo por categoría y de las órdenes por fecha & E3 & \textbf{Sin cambios.} La guía rectora lo dice literalmente: la persistencia distribuida de la E3 se mantiene & \reuso \\
		\addlinespace
		Prueba de tolerancia a fallos del clúster y consultas comparadas & E3 & Se conservan como evidencia y alimentan el panel de disponibilidad Raft & \reuso \\
		\addlinespace
		\emph{Pipeline} PySpark sobre el histórico de navegación y compras & E3 \newline PE-U4 & Se empaqueta \textbf{en el mismo \emph{compose}} y se conserva en \texttt{apps/spark/} & \cambia \\
		\addlinespace
		Métricas Prometheus del clúster (\texttt{crdb\_*}) & E3 & Se reutilizan y se les añaden cuatro métricas de aplicación & \cambia \\
		\addlinespace
		Análisis individual del modelo de consistencia y de la tecnología de datos & TA-IND-02 \newline FORM-IC-02 & Alimentan la sección de Fundamento teórico y Trabajos relacionados & \reuso \\
		\addlinespace
		\emph{Issues} abiertos por el equipo revisor & FORM-TL-05 & Debían responderse en la semana 17. Los que sigan abiertos hay que cerrarlos esta semana & \cambia \\
		\addlinespace
		Caso real documentado de las exposiciones & Cortes 1 y 2 & Sostiene la motivación del problema en la Introducción & \reuso \\
		\bottomrule
	\end{longtable}
\end{center}

\begin{avisoBox}{Lo que más se subestima de este inventario}
	El clúster de la E3 y el canal TCP de la PE-U1 llevan semanas sin ejecutarse. \textbf{Levántelos hoy}, antes que cualquier otra cosa. Si se rompieron por un cambio posterior, eso debía haberse detectado en la semana 15 y ahora consume días que no sobran: es lo primero que hay que resolver.
\end{avisoBox}

\section{Lo que la Entrega 4 pide de nuevo al equipo AGLS}

\subsection{Las dos aplicaciones cliente, que son el corazón de la entrega}

\begin{pisoBox}{Requisito estricto de la guía rectora}
	Cada equipo debe presentar dos aplicaciones cliente que consumen la misma API distribuida: una aplicación web y una aplicación móvil. Ambas son obligatorias y no sustituibles entre sí. \textbf{La ausencia total de una de ellas aplica un factor multiplicativo de 0,5 sobre la nota final de la entrega.}
\end{pisoBox}

La guía rectora concreta en su Tabla 1 qué debe ser cada una para AGLS:

\begin{center}
	\small
	\begin{tabular}{@{}p{2.5cm}p{12.5cm}@{}}
		\toprule
		\textbf{Web (SPA)} & Catálogo, carrito, \emph{checkout}, panel de vendedor y tablero de métricas para el administrador \\
		\addlinespace
		\textbf{Móvil} & Cliente móvil de compra: catálogo, carrito, notificaciones de estado del pedido y cámara para reconocer códigos de barras \\
		\addlinespace
		\textbf{Patrones GoF} & Repository, Factory Method, Strategy (métodos de pago), Observer (eventos de pedido) y Decorator (promociones) \\
		\bottomrule
	\end{tabular}
\end{center}

\begin{defBox}{Una coincidencia que conviene aprovechar}
	Los cinco patrones asignados a AGLS no son decorativos ni intercambiables: \textbf{Strategy} encapsula los métodos de pago, y es el mismo mecanismo con el que se conmuta la coordinación de la compra en la sección 5; \textbf{Factory Method} crea esas estrategias; \textbf{Observer} propaga los eventos del ciclo del pedido; \textbf{Decorator} aplica las promociones sobre el precio base; y \textbf{Repository} aísla el acceso al clúster. Los cinco están en el camino crítico del dominio, así que documentarlos en el ADR-005 es describir lo que el sistema hace, no adornarlo.
\end{defBox}

\subsection{Resumen de los ocho módulos}

\begin{center}
	\footnotesize
	\begin{longtable}{@{}p{1.0cm}p{3.5cm}p{1.6cm}p{9.0cm}@{}}
		\toprule
		\textbf{Mód.} & \textbf{Qué exige} & \textbf{Etiqueta} & \textbf{Qué significa concretamente para AGLS} \\
		\midrule
		\endhead
		A & Refactorización del \emph{backend} en capas & \cambia & El código de E2 y E3 se reorganiza en cuatro paquetes. La API externa no cambia \\
		\addlinespace
		B & Aplicación web (SPA) & \nuevo & Catálogo, carrito, \emph{checkout}, panel de vendedor y tablero del administrador \\
		\addlinespace
		C & Aplicación móvil & \nuevo & Cliente de compra con notificaciones de estado y lectura de códigos de barras \\
		\addlinespace
		D & Pirámide de pruebas & \cambia & Las pruebas de E2 y E3 se conservan y se les añaden contrato, extremo a extremo y carga \\
		\addlinespace
		E & \emph{Pipeline} de siete trabajos & \cambia & El flujo existente se amplía a siete trabajos y publica imágenes con el \emph{hash} del \emph{commit} \\
		\addlinespace
		F & Observabilidad distribuida & \cambia & Las métricas del clúster se conservan; se añaden cuatro de aplicación, trazas y panel de seis vistas \\
		\addlinespace
		G & Evaluación experimental ISO 25010 & \nuevo & El protocolo de medición completo. \textbf{Con la ampliación de la sección 5} \\
		\addlinespace
		H & Consolidación del documento & \cambia & El manuscrito acumulativo que integra E1 a E4 \\
		\bottomrule
	\end{longtable}
\end{center}

\section{Procedimiento por módulo}

\subsection{Módulo A --- Refactorización del \emph{backend} \emph{[Arquitecto]}}

\textbf{Entrada:} el microservicio principal de la E3 conectado al clúster CockroachDB. \textbf{Se conserva su comportamiento externo íntegro.}

\begin{enumerate}[leftmargin=*,itemsep=3pt]
	\item Reorganice el código de cada microservicio en cuatro paquetes: \texttt{presentation}, \texttt{application}, \texttt{domain} e \texttt{infrastructure}.
	\item Extraiga al paquete \texttt{domain} las entidades y objetos de valor del dominio: \texttt{Producto}, \texttt{Catalogo}, \texttt{Carrito}, \texttt{Orden}, \texttt{Pago}, \texttt{ReservaStock} y las reglas de compatibilidad. \textbf{Estas clases no llevan anotaciones de persistencia}: las anotaciones se quedan en las clases de \texttt{infrastructure}.
	\item Introduzca puertos, es decir interfaces, en \texttt{domain}, y sus adaptadores en \texttt{infrastructure}: repositorios sobre el clúster, clientes HTTP y publicadores de mensajes.
	\item Documente cinco patrones GoF en \texttt{docs/adr/ADR-005-\allowbreak patrones-gof.md} con el formato Nygard. Para AGLS: Repository, Factory Method, Strategy, Observer y Decorator. Cada uno debe nombrar las clases reales que ocupan cada rol.
	\item Actualice las pruebas de E2 y E3 para que se apoyen en dobles de prueba de los puertos, salvo las de integración.
\end{enumerate}

\textbf{Salida:} misma API externa, estructura interna desacoplada, cobertura unitaria $\geq 70\,\%$.

\begin{avisoBox}{La trampa más común de este módulo}
	Mover archivos de carpeta no es refactorizar en capas. La prueba es la dirección de las dependencias: si una clase de \texttt{domain} importa algo de \texttt{infrastructure}, la refactorización no está hecha. Ejecute esa comprobación antes de dar el módulo por cerrado.
\end{avisoBox}

\subsection{Módulo B --- Aplicación web \emph{[Desarrollador]}}

Todo el módulo es nuevo. Se crea en \texttt{apps/web/}.

\begin{enumerate}[leftmargin=*,itemsep=3pt]
	\item Proyecto con TypeScript en modo estricto y verificador de estilo obligatorio.
	\item Arquitectura por características: cada módulo de negocio agrupa componentes, servicios de acceso a la API, \emph{hooks} o servicios, y pruebas.
	\item Cubra al menos cinco rutas. Para AGLS: \texttt{/}, \texttt{/login}, \texttt{/catalogo}, \texttt{/settings} y \texttt{/about}, con el \emph{checkout} y el panel de vendedor tras rutas protegidas por credencial.
	\item El tablero del administrador debe mostrar las \textbf{métricas de venta y de estado de los pedidos}, que es lo que la guía rectora asigna a AGLS, alimentado por las métricas de negocio del Módulo F.
	\item Integre internacionalización en español e inglés, tema claro y oscuro, y estados explícitos de carga y de error.
	\item Consuma la API a través del API Gateway. Guarde la credencial en \emph{cookie} de solo servidor o en almacenamiento de sesión con expiración; \textbf{nunca en almacenamiento local plano}.
	\item Empaquete con \texttt{Dockerfile} multi-etapa, servido por Nginx, y agréguelo a \texttt{docker-compose.yml}.
	\item Publique el \texttt{.tar.gz} de la distribución como artefacto del \emph{pipeline}, versionado con el \emph{hash} del \emph{commit}.
\end{enumerate}

\subsection{Módulo C --- Aplicación móvil \emph{[Desarrollador]}}

Todo el módulo es nuevo. Se crea en \texttt{apps/mobile/}.

\begin{enumerate}[leftmargin=*,itemsep=3pt]
	\item Proyecto con nivel mínimo de interfaz de programación 26 y objetivo 34, o su equivalente si el equipo elige multiplataforma.
	\item Arquitectura de modelo, vista y modelo de vista, con un repositorio por característica.
	\item Autenticación contra el \textbf{mismo} \emph{endpoint} \texttt{/auth/login} del \emph{backend}; la credencial se guarda cifrada.
	\item Listado y detalle del recurso principal, que para AGLS es \textbf{el pedido del cliente}, con recarga por gesto y modo sin conexión mediante caché local.
	\item Integre al menos dos capacidades del dispositivo. Para AGLS la guía rectora asigna \textbf{cámara} (lectura de códigos de barras para añadir un producto al carrito) y \textbf{notificaciones} (cambio de estado del pedido). Ambas tienen sentido de dominio, así que no hay que forzarlas.
	\item Pruebas unitarias de los modelos de vista y al menos una prueba instrumentada de extremo a extremo.
	\item Genere el paquete de instalación firmado con clave de depuración y publíquelo como artefacto del \emph{pipeline} en \texttt{release/apk/}.
\end{enumerate}

\begin{defBox}{Por qué la app móvil no es un adorno en este proyecto}
	La notificación de cambio de estado del pedido es lo que cierra el lazo de la transacción distribuida: es donde el cliente se entera de que su compra se confirmó, se anuló o se compensó. Sin la app, el resultado de la coordinación se queda dentro del sistema y nunca llega al usuario, que es justo el punto que la evaluación del Módulo G quiere medir.
\end{defBox}

\subsection{Módulo D --- Pirámide de pruebas \emph{[Responsable de Calidad]}}

\begin{center}
	\footnotesize
	\begin{tabular}{@{}p{2.5cm}p{1.6cm}p{10.4cm}@{}}
		\toprule
		\textbf{Nivel} & \textbf{Etiqueta} & \textbf{Qué hacer} \\
		\midrule
		Unitarias & \cambia & Se conservan las de E2 y E3, adaptadas a los puertos. Cobertura $\geq 70\,\%$ en \emph{backend}, web y móvil \\
		\addlinespace
		Integración & \cambia & Pruebas de repositorio contra un CockroachDB efímero levantado por el propio marco de pruebas \\
		\addlinespace
		Contrato & \nuevo & Web y móvil como consumidores, \emph{backend} como proveedor. Es lo que garantiza que los dos clientes hablan la misma API \\
		\addlinespace
		Extremo a extremo & \nuevo & Suite para la web y suite instrumentada para la móvil \\
		\addlinespace
		Carga & \cambia & Dos escenarios obligatorios: \textbf{50 usuarios durante 5 minutos} y \textbf{rampa progresiva de 0 a 200 usuarios en 10 minutos} \\
		\bottomrule
	\end{tabular}
\end{center}

Cada tipo de prueba se ejecuta en un trabajo distinto del \emph{pipeline} y sus resultados se almacenan como artefactos.

\subsection{Módulo E --- \emph{Pipeline} de integración y entrega \emph{[Responsable de Calidad]}}

El archivo \texttt{.github/workflows/ci-cd.yml} define siete trabajos encadenados: \texttt{lint}, \texttt{test-backend}, \texttt{test-web}, \texttt{test-mobile}, \texttt{build-images}, \texttt{build-mobile-apk} e \texttt{integration}. Solo cuando todos están en verde se publican las imágenes etiquetadas con el \emph{hash} del \emph{commit}.

\begin{avisoBox}{Verificación que exige la rúbrica}
	El indicador de nivel 4 pide los siete trabajos \textbf{en verde en el último \emph{commit} de la entrega}. Compruébelo el día anterior, no el mismo día: un fallo en \texttt{build-mobile-apk} a última hora deja la Dimensión 5 en nivel bajo aunque todo lo demás funcione.
\end{avisoBox}

\subsection{Módulo F --- Observabilidad \emph{[Arquitecto]}}

\begin{enumerate}[leftmargin=*,itemsep=3pt]
	\item Añada el \emph{starter} de trazas al microservicio principal y configure el exportador hacia un colector agregado al \emph{compose}.
	\item \textbf{Reutilice las métricas del clúster de la E3} y añada cuatro de aplicación: total de peticiones etiquetado por ruta, método y código; duración de las peticiones; total de eventos de negocio; y sesiones activas.
	\item Emita registros en formato JSON en el \emph{backend} y en la móvil, incluyendo el identificador de traza para correlación.
	\item Construya el panel con las seis vistas obligatorias: tasa de peticiones, latencia en los percentiles 50, 95 y 99, tasa de errores de servidor, disponibilidad del consenso del clúster, uso de recursos por contenedor y latencia de extremo a extremo desde el cliente móvil.
	\item Exporte el panel como JSON versionado en \texttt{ops/grafana/\allowbreak pfc-dashboard.json}.
\end{enumerate}

\begin{defBox}{Qué evento de negocio debe contar AGLS}
	La métrica de eventos de negocio no es genérica. Para AGLS debe contar, como mínimo: órdenes creadas, órdenes confirmadas, órdenes canceladas o compensadas, y cobros rechazados por la pasarela. Esos cuatro contadores son los que después alimentan la evaluación del Módulo G.
\end{defBox}

\section{Ampliación del Módulo G para el equipo AGLS}

Esta es la única exigencia que esta guía añade a la guía rectora, y es pequeña en esfuerzo.

\subsection{Qué pide la guía rectora y qué le falta}

El Módulo G define el protocolo: sujeto de estudio, instrumentos, \textbf{bloque completo con diez repeticiones por escenario descartando la primera y la última}, métricas objetivas de la Tabla 2, análisis con media, desviación típica e intervalo de confianza al 95\,\%, y amenazas a la validez con al menos tres internas y dos externas.

Ese protocolo está completo salvo en un punto: tal como está, mide el sistema \emph{una sola vez} y contrasta cada métrica con su umbral. Eso responde a «¿cumple el sistema?», que es una evaluación de conformidad. No responde a «¿qué aporta la decisión de diseño que tomamos?», porque no hay nada contra lo que comparar. En TiendaTech esa decisión de diseño tiene nombre propio y la propia asignatura la señala como la que el equipo debe defender: confirmación en dos fases frente a saga con compensación.

\begin{uteqbox}{La ampliación en una frase}
	El mismo bloque, los mismos escenarios, las mismas diez repeticiones, ejecutados dos veces: con la transacción de compra coordinada por confirmación en dos fases y con ella coordinada por saga. Nada más cambia.
\end{uteqbox}

\subsection{Las tres adiciones, en orden de ejecución}

\subsubsection{Adición 1 --- Un condicional y una variable de entorno \emph{[Desarrollador]}}

En el punto donde hoy se confirma la compra, introduzca dos comportamientos seleccionables con la variable \texttt{COORD}:

\begin{center}
	\small
	\begin{tabular}{@{}p{1.4cm}p{4.0cm}p{9.6cm}@{}}
		\toprule
		\textbf{Valor} & \textbf{Comportamiento} & \textbf{Qué hace} \\
		\midrule
		\texttt{2pc} & Confirmación en dos fases & Un coordinador pregunta a Inventario, Pagos y Órdenes si pueden confirmar; solo si los tres responden afirmativamente, confirma. Mantiene bloqueos mientras dura la transacción \\
		\addlinespace
		\texttt{saga} & Saga con compensación & Cada servicio confirma su parte de inmediato; si un paso posterior falla se ejecutan las compensaciones en orden inverso: devolver stock, revertir cobro y cancelar la orden \\
		\bottomrule
	\end{tabular}
\end{center}

No son dos implementaciones ni dos ramas: es una estrategia seleccionable, que además es exactamente el patrón \textbf{Strategy} que la guía rectora ya le asigna a AGLS para los métodos de pago. La variable se declara en \texttt{.env.example}.

\subsubsection{Adición 2 --- El inyector de fallos de la pasarela y el oráculo \emph{[Desarrollador]}}

Un servicio intermediario entre el sistema y la pasarela de pago simulada, capaz de producir \textbf{omisión} (la petición se pierde y nunca llega respuesta) y \textbf{temporización} (la respuesta llega cinco segundos tarde, cuando el sistema ya pudo darla por perdida), con probabilidad y semilla configurables. Y un oráculo que, al terminar cada corrida, revisa la base de datos y responde si quedó todo cuadrado.

\begin{lstlisting}
experimentos/
  protocolo-e4.md            # ya exigido por la guia rectora
  inyector_pasarela.py       # NUEVO: omision y temporizacion con semilla
  oraculo_consistencia.py    # NUEVO: verifica los cuatro invariantes
  resultados/
    inconsistencias.csv      # NUEVO: ordenes con invariante violada, por corrida
    compatibilidad.csv       # NUEVO: matriz de confusion del asistente
    iso25010.csv             # ya exigido
    boxplot_latencia.png     # ya exigido
\end{lstlisting}

\begin{defBox}{Los cuatro invariantes que debe verificar el oráculo}
	\textbf{1.} Toda orden confirmada tiene un cobro registrado por el importe exacto de la orden. \textbf{2.} Toda orden confirmada descontó el stock del producto exactamente una vez. \textbf{3.} Ningún producto tiene stock negativo en ningún instante del registro. \textbf{4.} Toda orden cancelada o compensada devolvió el stock y no dejó cobro pendiente.\\[5pt]
	Sin este guion no existe la variable central de la ampliación, y sin ella el experimento no tiene resultado.
\end{defBox}

\subsubsection{Adición 3 --- Métricas nuevas en la tabla que ya existe \emph{[Responsable de Calidad]}}

\begin{center}
	\small
	\begin{tabular}{@{}p{4.2cm}p{6.2cm}p{4.6cm}@{}}
		\toprule
		\textbf{Métrica} & \textbf{Cómo se calcula} & \textbf{Qué responde} \\
		\midrule
		Tasa de inconsistencia observable & Proporción de órdenes que violan alguno de los cuatro invariantes al terminar la corrida & Cuánto cuesta en corrección cada estrategia \\
		\addlinespace
		Tasa de abortos & Proporción de compras iniciadas que no llegan a confirmarse & El coste comercial de la estrategia \\
		\addlinespace
		Convergencia de la compensación & Tiempo entre el fallo y la restitución completa del estado & Solo en la saga: cuánto dura el estado inconsistente \\
		\bottomrule
	\end{tabular}
\end{center}

\begin{avisoBox}{La inconsistencia y la latencia se leen juntas}
	La confirmación en dos fases tenderá a dar menos inconsistencia y más latencia; la saga, al revés. Reportar solo una de las dos convierte una decisión de ingeniería en una consigna. El valor del experimento está precisamente en poner las dos cifras una al lado de la otra.
\end{avisoBox}

\subsection{Cómo se analiza la comparación}

El Módulo G pide media, desviación típica e intervalo de confianza al 95\,\% para contrastar cada métrica con su umbral. Eso es lo correcto para una evaluación de conformidad. Para \emph{comparar dos configuraciones entre sí} hace falta algo más, y es barato de añadir porque son dos líneas en el cuaderno de análisis:

\begin{itemize}[leftmargin=*,itemsep=4pt]
	\item \textbf{Mediana además de media}, porque las latencias no se distribuyen de forma simétrica y la media se desplaza con unos pocos valores extremos.
	\item Una \textbf{prueba U de Mann-Whitney} entre configuraciones. Compara si un grupo tiende a dar valores mayores que el otro sin suponer que los datos siguen una distribución normal, que es justamente lo que no se cumple aquí.
	\item Un \textbf{tamaño del efecto} con el estadístico $\hat{A}_{12}$ de Vargha-Delaney, que se lee como la probabilidad de que una corrida tomada al azar de un grupo supere a una del otro. Responde a «¿la diferencia es grande?», que no es lo mismo que «¿la diferencia existe?».
	\item \textbf{Diagramas de caja} por escenario y configuración, no gráficos de barras con la media.
	\item Las proporciones se reportan con \textbf{intervalo de confianza binomial}, no como porcentaje desnudo.
\end{itemize}

\begin{avisoBox}{Un cero observado no es un cero garantizado}
	Si en las corridas no aparece ninguna inconsistencia bajo una estrategia, el resultado no es «nunca ocurre», sino «no se observó en $n$ corridas». El intervalo de confianza binomial da el rango honesto y debe escribirse así en el documento.
\end{avisoBox}

\subsection{Los escenarios del bloque experimental}

\begin{center}
	\small
	\begin{tabular}{@{}p{1.3cm}p{5.4cm}p{8.3cm}@{}}
		\toprule
		\textbf{Cód.} & \textbf{Escenario} & \textbf{Configuración} \\
		\midrule
		Esc-1 & Carga nominal & 50 usuarios durante 5 minutos, pasarela sin fallo. Es el escenario que la Tabla 2 usa para los umbrales \\
		\addlinespace
		Esc-2 & Compra concurrente con omisión & Rampa de 0 a 200 usuarios sobre el mismo producto, con la pasarela perdiendo el 10\,\% de las peticiones \\
		\addlinespace
		Esc-3 & Compra concurrente con temporización & Igual que el anterior, pero la pasarela responde cinco segundos tarde en el 10\,\% de los casos \\
		\addlinespace
		Esc-4 & Contención máxima sobre un producto & 200 usuarios comprando la última unidad disponible del mismo producto \\
		\bottomrule
	\end{tabular}
\end{center}

Los escenarios 2, 3 y 4 se ejecutan bajo \texttt{2pc} y bajo \texttt{saga}; el escenario 1 solo bajo la configuración que el equipo declare como predeterminada. Con las diez repeticiones del protocolo rector son \textbf{70 corridas}.

\begin{pisoBox}{Dos comprobaciones sin las cuales el experimento no vale}
	\textbf{Control negativo.} En el escenario 4, con la última unidad disponible y 200 usuarios simultáneos, el sistema debe registrar contención real: si todas las compras se resuelven sin conflicto, las peticiones no llegaron simultáneas y ninguna comparación posterior tiene sentido.\\[5pt]
	\textbf{Escenario 3 obligatorio.} La temporización es el modo de fallo que distingue de verdad a las dos estrategias, porque es el que deja al sistema sin saber si el cobro ocurrió. Sin él, el resultado favorece artificialmente a la saga.
\end{pisoBox}

\subsection{Amenazas a la validez}

La guía rectora pide al menos tres internas y dos externas. Para AGLS, las candidatas naturales son:

\begin{center}
	\small
	\begin{tabular}{@{}p{2.0cm}p{12.9cm}@{}}
		\toprule
		\textbf{Tipo} & \textbf{Amenaza y mitigación} \\
		\midrule
		Interna & Ruido de la máquina anfitriona y competencia entre contenedores. Se mitiga con el descarte de la primera y la última repetición \\
		Interna & Semilla del inyector de la pasarela. Se mitiga fijándola y declarándola por corrida \\
		Interna & Efecto del orden de ejecución de las dos estrategias. Se mitiga alternando el orden entre bloques \\
		\addlinespace
		Externa & El catálogo y el patrón de compra son sintéticos, y la pasarela está simulada \\
		Externa & La escala del despliegue local frente a la de un comercio en operación real \\
		\bottomrule
	\end{tabular}
\end{center}

\subsection{Qué debe quedar archivado, y por qué no se puede reconstruir después}

Las mediciones de latencia y rendimiento no significan nada sin el entorno en que se tomaron. Esa información se pierde en cuanto el equipo actualiza una dependencia o cambia de máquina, y \textbf{no hay forma de recuperarla más tarde}. Debe quedar registrada en \texttt{experimentos/protocolo-e4.md} el mismo día de las corridas.

\begin{center}
	\small
	\begin{tabular}{@{}p{4.4cm}p{10.5cm}@{}}
		\toprule
		\textbf{Qué se archiva} & \textbf{Detalle exigido} \\
		\midrule
		Entorno de ejecución & Modelo de procesador, número de núcleos, memoria RAM, tipo de disco, sistema operativo y versión, versión del motor de contenedores \\
		\addlinespace
		Versiones exactas & Etiqueta o \emph{hash} de cada imagen usada, versión del entorno de ejecución del \emph{backend}, del clúster y de la herramienta de carga \\
		\addlinespace
		Semillas & Semilla del generador de compras y del inyector de la pasarela, una por corrida \\
		\addlinespace
		Datos crudos & Los archivos de \textbf{todas} las corridas, no solo los resúmenes. Los resúmenes se pueden recalcular; los crudos, no \\
		\addlinespace
		Informe del oráculo & El resultado de los cuatro invariantes en cada corrida, emparejado con su identificador \\
		\addlinespace
		Cuaderno de análisis & El que regenera todas las figuras y tablas a partir de los crudos, sin intervención manual \\
		\addlinespace
		Instantánea citable & Una versión congelada del repositorio con identificador persistente, generada al cerrar la entrega \\
		\bottomrule
	\end{tabular}
\end{center}

\begin{defBox}{La prueba de que el archivo está completo}
	Otra persona, con solo el repositorio y el \texttt{README.md}, debe poder clonar, levantar el sistema, ejecutar una corrida y \textbf{regenerar todas las figuras del documento} sin preguntar nada al equipo. Si necesita preguntar algo, eso que falta es precisamente lo que hay que archivar.
\end{defBox}


\section{Estructura del repositorio: qué existe y qué nace}

La guía rectora fija la estructura final. Esta tabla indica el origen de cada carpeta para que el equipo sepa qué crear y qué solo mover.

\begin{center}
	\footnotesize
	\begin{longtable}{@{}p{5.4cm}p{1.9cm}p{7.5cm}@{}}
		\toprule
		\textbf{Ruta} & \textbf{Etiqueta} & \textbf{Observación para AGLS} \\
		\midrule
		\endhead
		\texttt{docs/adr/ADR-001..004} & \reuso & Se conservan; se les añade el impacto de la E4 \\
		\texttt{docs/adr/ADR-005-\allowbreak patrones-gof.md} & \nuevo & Los cinco patrones asignados a AGLS, formato Nygard \\
		\texttt{docs/adr/ADR-006-\allowbreak eleccion-movil.md} & \nuevo & Justificación con al menos dos criterios cuantitativos \\
		\texttt{docs/api/openapi.yaml} & \cambia & Verificar que describe la API real tras el Módulo A \\
		\texttt{docs/diagrams/c4-nivel1,2} & \reuso & Sin cambios \\
		\texttt{docs/diagrams/c4-nivel3} & \cambia & Se actualiza con la web y la móvil \\
		\texttt{db/schema.sql} & \reuso & El esquema de la E3 \\
		\texttt{services/} & \cambia & Refactorización en cuatro capas \\
		\texttt{apps/web/} & \nuevo & Módulo B \\
		\texttt{apps/mobile/} & \nuevo & Módulo C \\
		\texttt{apps/spark/} & \cambia & Heredado de la E3; ahora dentro del mismo \emph{compose} \\
		\texttt{experimentos/\allowbreak protocolo-e4.md} & \nuevo & El protocolo del Módulo G ampliado \\
		\texttt{experimentos/\allowbreak inyector\_pasarela.py} & \nuevo & Adición 2 de esta guía \\
		\texttt{experimentos/resultados/} & \nuevo & Datos crudos, incluida la verdad de campo \\
		\texttt{ops/grafana/\allowbreak pfc-dashboard.json} & \nuevo & Panel exportado como código \\
		\texttt{ops/prometheus/} & \cambia & Se añaden las cuatro métricas de aplicación \\
		\texttt{ops/otel-collector/} & \nuevo & Colector de trazas \\
		\texttt{tests/contract/} & \nuevo & Contratos de los dos clientes \\
		\texttt{tests/e2e-web/} & \nuevo & Suite de la web \\
		\texttt{tests/integration/},\allowbreak{} \texttt{tests/load/} & \cambia & Se conservan y se amplían con los dos escenarios de carga \\
		\texttt{release/apk/},\allowbreak{} \texttt{release/screenshots/} & \nuevo & Artefactos generados por el \emph{pipeline} \\
		\texttt{.env.example} & \cambia & Añadir las variables de web, móvil y \texttt{COORD} \\
		\texttt{docker-compose.yml} & \cambia & \emph{Backend}, clúster, web y colector \\
		\bottomrule
	\end{longtable}
\end{center}

\section{Del material acumulado al documento final}

El manuscrito de la E4 es acumulativo: integra actualizadas todas las secciones de E1, E2 y E3, y las cita internamente, nunca como fuentes externas. Mínimo 20 páginas de contenido y 12 referencias en formato IEEE con identificador verificable.

\begin{center}
	\footnotesize
	\begin{longtable}{@{}p{3.4cm}p{1.0cm}p{2.6cm}p{7.6cm}@{}}
		\toprule
		\textbf{Sección} & \textbf{Págs.} & \textbf{Procede de} & \textbf{Qué hay que hacer con ese material} \\
		\midrule
		\endhead
		Resumen bilingüe & 1 & Nuevo & Máximo 200 palabras cada versión \\
		Introducción & 1,5 & E1 y exposiciones & Verificar las cifras y las fuentes; añadir las contribuciones del semestre \\
		Trabajos relacionados & 2 & TA-IND-02, FORM-IC-02 & Ampliar a un mínimo de seis fuentes indexadas sobre bases distribuidas, microservicios, móvil y calidad \\
		Fundamento teórico & 3 & E1 y E3 & Capas y SOLID, patrones, móvil, integración continua, observabilidad e ISO 25010 \\
		Arquitectura & 2 & E1 y E2 & C4 de nivel 1 a 3, los seis ADR y el \emph{stack} consolidado \\
		Aplicación web & 2 & Nuevo & Manual del Módulo B, capturas de las cinco rutas y evidencia de internacionalización \\
		Aplicación móvil & 2 & Nuevo & Manual del Módulo C, capturas y evidencia de las dos capacidades usadas \\
		Pruebas y CI/CD & 2 & E2, E3 y nuevo & Pirámide, resumen del \emph{pipeline} y cobertura \\
		Observabilidad & 1,5 & E3 y nuevo & Panel como figura y ejemplo de traza distribuida \\
		Evaluación ISO 25010 & 2 & Nuevo & Tabla de métricas con intervalo de confianza y discusión, \textbf{incluida la comparación de la sección 5} \\
		Discusión y amenazas & 1,5 & Nuevo & Interpretación, alcance y limitaciones \\
		Conclusiones & 1 & Todo & Cierre del ciclo E1 a E4 \\
		\bottomrule
	\end{longtable}
\end{center}

\begin{avisoBox}{Dónde entra la ampliación del Módulo G en el documento}
	No crea una sección nueva. La comparación entre configuraciones entra en \textbf{Evaluación ISO 25010} como un bloque de resultados con su tabla, y su lectura entra en \textbf{Discusión y amenazas}. La estructura de la Tabla 3 de la guía rectora se respeta sin tocar.
\end{avisoBox}

\section{Lo que debía estar terminado}

Antes del plan de los días que quedan, conviene fijar con precisión qué exigían las semanas ya transcurridas. Esta tabla no es un reproche: es el mapa para saber dónde mirar si algo falla.

\begin{center}
	\small
	\begin{tabular}{@{}p{1.8cm}p{5.8cm}p{7.3cm}@{}}
		\toprule
		\textbf{Semana} & \textbf{Qué debía quedar terminado} & \textbf{Cómo comprobarlo hoy en un minuto} \\
		\midrule
		4 & E1 y PE-U1: análisis, C4, ADR-001, sockets, gRPC y relojes de Lamport & El servidor de sockets arranca y los CSV de latencia están en el repositorio \\
		\addlinespace
		7 & E2: microservicios, puerta de enlace, OpenAPI, autenticación y \emph{compose} & \texttt{docker compose up} levanta y \texttt{/auth/login} devuelve credencial \\
		\addlinespace
		11 & E3: clúster de tres nodos, fragmentación y prueba de tolerancia a fallos & Los tres nodos responden y existe el vídeo de la prueba \\
		\addlinespace
		14 & \emph{Pipeline} PySpark y medición de \emph{speedup} & Los resultados de tiempos están en el repositorio \\
		\addlinespace
		16 & Revisión cruzada: \emph{issues} recibidos del equipo revisor & Ningún \emph{issue} abierto sin respuesta \\
		\addlinespace
		17 & Módulos A, B y C encaminados; el grueso del código nuevo & La web y la móvil compilan, aunque estén incompletas \\
		\bottomrule
	\end{tabular}
\end{center}

\section{Plan de los nueve días restantes}

El cierre es el viernes 4 de septiembre a las 23:55. Este plan supone que los Módulos A, B y C están al menos empezados. Si no lo están, aplique primero el triaje de la sección 10.

\begin{center}
	\small
	\begin{longtable}{@{}p{2.6cm}p{1.5cm}p{5.2cm}p{5.0cm}@{}}
		\toprule
		\textbf{Día} & \textbf{Semana} & \textbf{Trabajo} & \textbf{Al acostarse debe estar hecho} \\
		\midrule
		\endhead
		Jueves 27 & 17 & Auditoría de la sección 2. Levantar clúster y canal TCP. Cerrar \emph{issues} de la revisión cruzada & El inventario marcado y el sistema heredado en pie \\
		\addlinespace
		Viernes 28 & 17 & Módulo A: dependencias hacia el dominio, ADR-005 con los cinco patrones & Ninguna clase de \texttt{domain} importa \texttt{infrastructure} \\
		\addlinespace
		Sábado 29 & 17 & \textbf{Módulo C, día completo}: autenticación, listado y detalle & La app móvil arranca, autentica y lista pedidos reales \\
		\addlinespace
		Domingo 30 & 17 & Módulo C: lectura de códigos de barras y notificaciones. Módulo B: rutas y \emph{login} & Las dos capacidades del dispositivo funcionan \\
		\addlinespace
		Lunes 31 & 18 & Módulo B: catálogo, carrito, \emph{checkout} y tablero, idiomas, tema, estados & Las cinco rutas navegables y protegidas \\
		\addlinespace
		Martes 1 & 18 & Adiciones 1 y 2: \texttt{COORD} y el inyector de la pasarela. Módulo F & \texttt{CORREL} conmuta; el panel muestra las seis vistas \\
		\addlinespace
		Miércoles 2 & 18 & \textbf{Módulo G}: las 100 corridas, lanzadas por lotes desde la mañana & Los datos crudos y los informes del oráculo en el repositorio \\
		\addlinespace
		Jueves 3 & 18 & Módulos D y E: pruebas de contrato y extremo a extremo; los siete trabajos en verde. Análisis y figuras & El \emph{pipeline} verde y todas las figuras generadas \\
		\addlinespace
		Viernes 4 & 18 & Módulo H: manuscrito, fusión a \texttt{main} por \emph{pull request}, verificación final & \textbf{Último \emph{commit} antes de las 23:55} \\
		\bottomrule
	\end{longtable}
\end{center}

\begin{avisoBox}{Dos decisiones de calendario que conviene no discutir}
	\textbf{El sábado es para la aplicación móvil.} Es el módulo con más sorpresas: firma del paquete, permisos del dispositivo, emulador, versiones del entorno de compilación. Dejarla para el final es la forma más común de perder la mitad de la nota.\\[5pt]
	\textbf{Las corridas del Módulo G se lanzan el miércoles por la mañana.} Cien corridas por lotes no caben en una noche, y si algo falla a mitad hace falta un día de margen para repetirlas.
\end{avisoBox}

\section{Triaje: qué hacer si algo no está}

Con nueve días no cabe todo. Si el equipo llega tarde, la peor decisión es repartir el esfuerzo por igual. Esta tabla ordena el trabajo por \textbf{lo que cuesta no tenerlo}, de mayor a menor.

\begin{center}
	\small
	\begin{longtable}{@{}p{0.9cm}p{5.4cm}p{3.0cm}p{5.0cm}@{}}
		\toprule
		\textbf{Ord.} & \textbf{Si falta esto} & \textbf{Coste} & \textbf{Mínimo aceptable en un día} \\
		\midrule
		\endhead
		1 & La aplicación móvil, por completo & \textbf{Factor 0,5 sobre toda la nota} & Que arranque, autentique contra \texttt{/auth/login}, liste pedidos y use dos capacidades del dispositivo \\
		\addlinespace
		2 & La aplicación web, por completo & \textbf{Factor 0,5 sobre toda la nota} & Cinco rutas, \emph{login} con credencial y el catálogo con carrito y \emph{checkout} \\
		\addlinespace
		3 & Los datos del Módulo G & Hunde la evaluación y la discusión & Reducir a los escenarios 1, 3 y 4, manteniendo las diez repeticiones \\
		\addlinespace
		4 & El \emph{pipeline} con los siete trabajos & Dimensión 5 en nivel bajo & Que los siete existan y estén en verde, aunque las pruebas sean mínimas \\
		\addlinespace
		5 & El manuscrito de 20 páginas & Dimensión 8 completa & Escribir sobre lo medido, nunca sobre lo planeado \\
		\addlinespace
		6 & La refactorización en capas & Dimensión 1 & Al menos el microservicio principal en cuatro paquetes \\
		\addlinespace
		7 & El panel con las seis vistas & Dimensión 6 & Las seis vistas, aunque sin refinar \\
		\bottomrule
	\end{longtable}
\end{center}

\begin{pisoBox}{Lo que nunca se recorta, aunque falte tiempo}
	\textbf{El control negativo} y \textbf{el escenario de dos averías simultáneas}. Sin el primero no se sabe si la carga reproduce el fenómeno; sin el segundo, la fusión errónea queda oculta y el resultado sale bueno y falso. Antes que recortarlos, reduzca el número de escenarios.\\[5pt]
	\textbf{El registro del entorno de ejecución.} Son diez minutos de trabajo y es el único dato de toda la entrega que no se puede reconstruir después.
\end{pisoBox}

\begin{avisoBox}{Sobre reducir las repeticiones}
	La tentación evidente es bajar de diez repeticiones a tres. \textbf{No lo haga.} Es preferible ejecutar menos escenarios con las diez repeticiones completas que todos los escenarios con tres: con tres corridas, el intervalo de confianza sale tan ancho que ninguna comparación resulta concluyente y el trabajo de medir se pierde entero.
\end{avisoBox}

\section{Defensa oral}

Además de las ocho preguntas obligatorias del Anexo B de la guía rectora, el equipo AGLS debe poder responder estas cuatro:

\begin{enumerate}[leftmargin=*,itemsep=3pt]
	\item ¿Cuál fue la tasa de inconsistencia observable con cada estrategia de coordinación, con su intervalo de confianza?
	\item Si la pasarela falla justo después de descontar el stock, ¿qué ocurre con cada estrategia y en cuánto tiempo se resuelve el estado?
	\item ¿Qué precio en latencia paga la estrategia más consistente, y qué significa ese número para un cliente que está pagando?
	\item ¿Dónde se ubica TiendaTech en el teorema CAP y por qué la respuesta es distinta para el inventario y para el catálogo?
\end{enumerate}

\section{Lista de verificación final}

A la lista del Anexo A de la guía rectora, el equipo AGLS añade estas comprobaciones:

\begin{center}
	\small
	\begin{tabular}{@{}p{0.7cm}p{14.2cm}@{}}
		\toprule
		\textbf{$\square$} & \textbf{Verificación adicional para AGLS} \\
		\midrule
		$\square$ & \texttt{COORD} conmuta entre \texttt{2pc} y \texttt{saga} sin tocar código ni cambiar de rama \\
		$\square$ & El inyector produce omisión y temporización de forma reproducible con la misma semilla \\
		$\square$ & \textbf{Control negativo validado}: con 200 usuarios sobre la última unidad hay contención real \\
		$\square$ & El escenario de temporización está ejecutado, no solo el de omisión \\
		$\square$ & Inconsistencia y latencia se reportan juntas, nunca una sola \\
		$\square$ & Las diez repeticiones por escenario están completas y se descartaron la primera y la última \\
		$\square$ & El tablero del administrador muestra métricas de venta reales en la consola web \\
		$\square$ & La app móvil lee un código de barras y recibe la notificación de estado del pedido \\
		$\square$ & Los cuatro contadores de eventos de negocio de AGLS aparecen en el panel \\
		$\square$ & Los cinco patrones del ADR-005 se corresponden con clases reales del código \\
		$\square$ & El entorno de ejecución está registrado: procesador, núcleos, RAM, disco, sistema y versiones \\
		$\square$ & Están los datos crudos de las 70 corridas, no solo los resúmenes \\
		$\square$ & La comparación entre estrategias reporta mediana, prueba estadística y tamaño del efecto \\
		$\square$ & El cuaderno de análisis regenera todas las figuras sin intervención manual \\
		\bottomrule
	\end{tabular}
\end{center}

\vspace{6pt}
\begin{uteqbox}{Nota final para el equipo}
	La Entrega 4 parece enorme porque son ocho módulos, pero la mayor parte del sistema ya está construido: el clúster, los microservicios, la puerta de enlace y el \emph{pipeline} de datos vienen de las entregas anteriores y se conservan. Lo verdaderamente nuevo son las dos aplicaciones cliente y el instrumental de medición. Y la ampliación de la sección 5 cabe en un condicional, dos guiones cortos y tres filas más en una tabla que de todos modos hay que llenar. La decisión entre confirmación en dos fases y saga lleva todo el semestre sobre la mesa: esta es la semana de resolverla con datos propios en lugar de con argumentos.
\end{uteqbox}

\end{document}
